\section{Domande al prof}
\begin{enumerate}
    \item Io sto pensando alla simulazione come costruita a due livelli: il primo  si occupa della generazione degli ordini e dell'ottimizzazione periodica, mentre il secondoe simula gli spostamenti dei pod. E questi due livelli devono in qualche modo interagire in quanto (1) lo stato attuale del sistema influenza alcuni vincoli dell'ottimizzatore (ad esempio se non ho finito di servire l'ordine $1$ alla workstation $2$ devo inserire un vincolo per cui $x_1,2,0 = 1$), e viceversa (2) l'output dell'ottimizzatore è l'input della simulazione dei movimenti. \\
    Per farli interagire senza incoerenze, posso assumere che quando parte lottimizzazione il sistema sia \enquote{congelato} e poi quando faccio ripartire la simulazione resetto la lista degli eventi? \\
    Alternativamente, potrei assumere (mi sembra più drastico) che l'ottimizzatore parta solo quando ho finito di servire tutti gli oridini considerati nell'ottimizzazione precedente.

    \item Se considero le variabili $y_{pwt}$ come valide indicazioni di scheduling dei pod, in quale istante della simulazione devo effettivamente attivare la richiesta di movimentazione del pod verso la workstation? \\
    Mi vengono in mente due approcci: (1) impostare una coda delle richieste che attinge ai robot come risorsa (problema: come decido chi ha la priorità tra pod necessari allo stesso step) oppure (2) generare richieste \enquote{anticipate} che tengano conto di quanto disti il pod dalla workstation. \\
    Seconda questione: Se ho dei robot liberi può avere senso che inizio a prelevare pod che mi serviranno?

    


\end{enumerate}


\section{Questioni a cui devo pensare}
\begin{enumerate}
    \item PROBLEMA: fare attenzione ai casi $y_{p,w_1,t} = 1$ e $y_{p,w_2,t+1} = 1$ $\longrightarrow$  forse risolto
    \item PROBLEMA: se $y_{p,w_1,t} = 1$ e $y_{p,w_2,t+2} = 1$ ha effettivamente senso che il pod torni al suo posto? C'è qualcosa nell'ottimizzatore che si assicuri che ciò non accada?
\end{enumerate}
